% !PRINTED 581
\weberpage{581}
%<<CONT
oder endlich nach (2) und (3)
\weq{12}{-27\gamma = nA + 3n - 1.}

Aus dieser Relation folgt, da $\gamma$ eine ganze Zahl ist,
\weq{13}{A \equiv 1\ (\mathrm{mod}\ 3)}
und daraus $A^3 \equiv 1\ (\mathrm{mod}\ 9)$. Nun ist nach (2), (3) und (6)
\weq{14}{n = a^2 - ab + b^2}
oder
\weq{15}{4n = A^2 + 3b^2}
und wenn wir diesen Werth in (12) substituiren,
\[
-4\cdot 27\gamma = A^3 + 3Ab^2 + 3A^2 + 9b^2 - 4,
\]
und daraus
\[
3b^2 \equiv 0\ (\mathrm{mod}\ 9).
\]

Es ist also $b$ durch 3 theilbar, und wenn wir $b = 3B$ setzen,
so folgt aus (15)
\weq{16}{4n = A^2 + 27B^2.}

Wir können auch leicht die Quadratwurzel aus der Dis\-criminante bilden, nämlich
\weq{17}{\sqrt{D} = (\eta - \eta_1)(\eta - \eta_2)(\eta_1 - \eta_2) = s - s',}
wenn wir (10) und (11) von einander subtrahiren. Man erhält
daraus
\weq{18}{\sqrt{D} = nB.}

Die cubische Gleichung, deren Wurzeln die drei Grössen $\eta$
sind, hat also, was übrigens schon aus dem am Anfang von
\S.\,171 Bemerkten folgt, drei reelle Wurzeln. Sie vereinfacht
sich, wenn man
\[
3\eta + 1 = \xi
\]
setzt, und ergiebt dann
\weq{19}{\xi^3 - 3n\xi - nA = 0.}

Substituirt man auf der linken Seite für $\xi$ der Reihe nach
\weq{20}{-2\sqrt{n}, \quad -\sqrt{n}, \quad +\sqrt{n}, \quad +2\sqrt{n},}
so ergeben sich die Werthe
\[
-n(2\sqrt{n} + A), \quad +n(2\sqrt{n} - A), \quad -n(2\sqrt{n} + A),
\]
\[
+n(2\sqrt{n} - A),
\]
die, weil nach (16) $A$ absolut kleiner als $2\sqrt{n}$ ist, abwechselnde
Vorzeichen haben. Es liegt also in jedem der drei durch die
Werthe (20) begrenzten Intervalle \gesperrt{eine} Wurzel der Gleichung (19).
